% Page 086

\seite{86}

und {[trafen]} gegen 8 Uhr {[morgens]}, von [Maschinen/Männern?]\ob
\unsicher\unsicher{} [geschmückt/geräumt]\unsicher{}\unsicher\unsicher{} {[Bestimmungs]}ort\ob
ein, während {[in]}\unsicher{}\unsicher Zeit die alte Glocke, die\ob
anno 1793 in Masburg\unsicher{}\unsicher {[der]} [Praxis?]\unsicher{}\unsicher\ob
[für/von]\unsicher{}\unsicher\unsicher{} in\unsicher{}\unsicher von der {[Glockengießerei]}\unsicher{} in\ob
\unsicher\unsicher{} [gegossen/gezogen]\unsicher{}\unsicher [Abholung?] [fertig?] {[war]}.

\margmark{5. Juni}%
Der Unterricht hat wieder begonnen.\ob
Die neuen Glocken wurden {[in]} der Pfingstwoche\ob
in der {[Kapelle]} auf zwei großen [Holzböcke/Gestelle] auf=\ob
gehängt zur Einweihung. Die größere von {[den]} {[beiden]}\ob
\margmark{Glocken}%
Glocken {[wiegt]}\unsicher{} Pfund, ist\unsicher{} abgestimmt, trägt\ob
das Bild des Gekreuzigten und von Johannes\ob
\unsicher [Beschriftung/Inschrift] {[mit]}\unsicher. [Hl.] Johannes\ob
{[befindet]} sich {[an]} {[der]} {[Gottes]}mutter [Seite?]\unsicher. {[Die]} La=\ob
\unsicher\unsicher{}\unsicher die\unsicher{}\unsicher ist:\unsicher{}\unsicher\unsicher.

Die kleinere Glocke trägt ebenfalls das Bild\ob
der Gekreuzigten, sowie von der Maria mit dem {[Jesus]}=\ob
kind. Diese Glocke {[wiegt]}\unsicher{} Pfund, {[ist]} {[ebenfalls]} abge=\ob
stimmt und trägt den {[Spruch]}:\unsicher{} [mich/ruft]\unsicher{} in allen\ob
{[Stund]} — das {[Wort]} {[auf]}\unsicher{} [zu/an] {[den]}\unsicher{} {[Mund]}.\unsicher{}\unsicher\unsicher.

Die [Rechnungen/Benennung] [beider?]\unsicher{} die [Lieferungen/Erzeugnisse] {[der]}\ob
Firma Mabilon [\&] Co in Saarburg\unsicher{} {[die]} hat vor=\ob
\unsicher die Glocken [gegossen/geliefert] {[worden]}.

Die Einweihung der Glocken fand {[heute]} nach=\ob
mittags {[um]} 3 Uhr statt, {[an]} der sich die ganze Gemeinde\ob
{[beteiligte]}. [Von/Die] auswärtigen Geistlichen\unsicher{} {[teil]}=\ob
nahmen\unsicher{} die Pfarrer von Malberg,\unsicher, Neuerburg,\ob
\unsicher von\unsicher{} {[Haupt]} Pfarrer [Krämer?]\unsicher\ob
\unsicher {[eine]} zusammen{[gestellte]} geistliche\unsicher{} {[Feier]}\unsicher{} und\ob
{[unter]} {[der]} Leitung des Lehrers Philipp {[sangen]} die {[Schul]}=\ob
{[kinder]} [ihre/die] Lieder: [\enquote{]Schwinge dich Glocke[...}]\ob
und [\enquote{]Die Glocken\unsicher{} ruft und grüß{[t]} die {[Frommen]}[}].

Im Anschluß an die Feier folgte die übliche {[Glocken]}=\ob
\unsicher [Schlag?] mit [Bändern/Girlanden] an die beiden Glocken\ob
\unsicher [der/die] Anwesenden {[eine]} [glänzende/glückliche]\unsicher{} [Verfassung?]\ob
\unsicher eines wohlgemeinten [Segenswunsches/Festaktes].
\pend
